\chapter{General Conclusion and Future Perspectives}
\label{chp:conclusion}

This thesis has undertaken a comprehensive investigation into the design and application of advanced fault-tolerant control strategies for autonomous aerial vehicles. Motivated by the increasing demand for reliability and safety in unmanned systems, the primary objective was to develop robust control architectures capable of guaranteeing stability and precise trajectory tracking even in the presence of severe system uncertainties, external disturbances, and actuator and/or sensor faults. Through the integration of nonlinear control theory, intelligent approximation techniques, and metaheuristic optimization, this work has successfully addressed the limitations of conventional control methods, offering rigorous solutions for distinct classes of aerial platforms ranging from grounded prototypes to complex multi-agent swarms.

The core theoretical contribution of this dissertation lies in the successful application of fixed-time stability concepts across different aerial platforms. Unlike finite-time control, where the convergence time depends on the initial states of the system, the fixed-time control laws developed herein ensure that tracking errors converge to a small neighborhood of the origin within a pre-calculated settling time that is independent of initial conditions. This property is paramount for safety-critical missions where time constraints are stringent. 

For Twin-Rotor UAV systems, a robust fault-tolerant control scheme was established by combining backstepping technique and Radial Basis Function Neural Networks (RBFNN). The inclusion of the Grey Wolf Optimizer (GWO) to tune the control parameters eliminated the subjectivity of manual tuning, resulting in a system that effectively handles the strong coupling and gyroscopic effects inherent to the Twin-Rotor system while compensating for sensor faults.

Building upon the single-agent results, the research extended to the more agile and underactuated Quadrotor UAV platform. Here, the complexity was managed through an Adaptive Fuzzy Backstepping controller. By leveraging the universal approximation capabilities of Fuzzy Logic Systems, the controller was able to estimate and mitigate the effects of aerodynamic disturbances and time-varying actuator faults online. The optimization framework was further refined by employing the Whale Optimization Algorithm (WOA), which demonstrated superior exploration capabilities in high-dimensional parameter spaces compared to other metaheuristics. The simulation results confirmed that this approach minimizes mechanical stress on the actuators while maintaining high-precision tracking, thereby enhancing the operational longevity of the UAV.

The final and most complex contribution of this work addressed the distributed formation control of Coaxial UAV swarms. Recognizing the limitations of static neural network structures in dynamic environments, a novel Self-Structuring Neural Network (SSNN) was introduced. This intelligent approximator dynamically adjusts its structure by adding or pruning neurons based on the complexity of the uncertainties and faults encountered during the mission. Combined with a distributed fixed-time protocol, this approach ensured that the swarm maintained its geometric formation and synchronized flight despite communication topology switches and the partial failure of individual agents. This represents a significant step forward in the autonomy of multi-agent systems, reducing the computational burden on individual agents while maximizing collective robustness.

The contributions of this thesis can be summarized along three complementary axes that collectively advance the state of the art in fault-tolerant control of autonomous aerial vehicles. First, on the theoretical front, this work establishes a unified fixed-time stability framework that guarantees convergence of tracking errors within a pre-calculated settling time independent of initial conditions, thereby overcoming a fundamental limitation of conventional finite-time approaches in safety-critical missions. Second, the research extends the scope from single-agent control to multi-UAV systems, progressing from a Twin-Rotor platform stabilized through a backstepping scheme augmented with Radial Basis Function Neural Networks and tuned via the Grey Wolf Optimizer, to a Quadrotor UAV governed by an Adaptive Fuzzy Backstepping controller optimized through the Whale Optimization Algorithm, and finally to a distributed formation protocol for Coaxial UAV swarms driven by a novel Self-Structuring Neural Network capable of dynamically adjusting its topology in response to uncertainties and partial agent failures. Third, on the practical front, the integration of metaheuristic optimization with intelligent approximators removes the subjectivity inherent to manual gain tuning, reduces actuator stress, and ensures synchronized formation flight under switching communication topologies, thus offering a coherent and scalable solution that bridges single-agent robustness and multi-agent cooperative resilience.

\section*{Future Perspectives}

While the control strategies proposed in this thesis have demonstrated robust performance and theoretical soundness, several avenues for future research remain open to further enhance the capabilities and real-world applicability of autonomous aerial systems. The most immediate progression of this work would be the validation of the proposed algorithms on experimental hardware platforms. Implementing the Fixed-Time Fault-Tolerant controllers on embedded systems, such as Pixhawk or dSPACE boards, would provide valuable insights into the effects of measurement noise, communication delays, and computational constraints that are often idealized in simulation environments. Real-time implementation would also necessitate the optimization of the code efficiency, particularly for the Self-Structuring Neural Networks, to ensure they can operate within the limited processing power of onboard microcontrollers.

Furthermore, the scope of fault scenarios considered in this work could be expanded to include more catastrophic and complex failure modes. Future research could investigate the system's response to structural damage, such as the partial loss of a propeller or frame deformation, which alters the fundamental aerodynamics of the vehicle beyond simple gain loss. Additionally, the resilience of the control systems against cyber-physical attacks, such as sensor spoofing or denial-of-service attacks on the inter-agent communication links in swarms, is a critical area that warrants investigation. Integrating cyber-security measures with fault-tolerant control would result in a truly resilient cyber-physical system.

Finally, the concepts developed for homogeneous coaxial swarms could be extended to heterogeneous multi-agent systems. A cooperative framework involving different types of unmanned systems, such as the collaboration between fixed-wing UAVs for long-range surveillance and rotary-wing UAVs for precise hovering, or even coordination between aerial and ground vehicles (UAV-UGV teams), presents a challenging control problem. Developing distributed fixed-time controllers that can account for the differing dynamic constraints and capabilities of heterogeneous agents would significantly broaden the operational scope of autonomous missions in search and rescue or environmental monitoring applications.